

\illustration{0.3}{tourist.png}{
	Generated using \href{https://www.canva.com}{canva.com} with prompt
	``Traveller with backpack and facemask, in front of mountains''.
}

% optionally define variables/limits for this problem






Long, long ago on a planet far, far away, a highly contagious virus caused an enduring pandemic.

Even so, the people wanted to travel between countries for their summer holidays.
In the good old before-days, travelling from any country to any other country took 1 full day.
However, during the pandemic,
certain countries preferred not to receive travellers from areas that had higher infection rates,
so they made them quarantine for a certain number of days before allowing them to continue their trip or start their holiday.

To keep everything fair, an independent Bureau for Accurate Pandemic Classification was founded. They assigned a $r$-value to each country based on the infection rate in that country. A higher $r$-value indicates higher infection rate.

Each country asked tourists to quarantine if the country they just came from had a $r$-value significantly higher than their own. In particular, when you wanted to travel from country $i$ to country $j$, you would have to quarantine for $t_j$ days if $r_i > r_j + m$.

Archaeologists have found evidence of $q$ tourists travelling between $n$
countries. For each tourist, the start and destination are known. The question
that remains to be answered is: how long was each tourist's minimal travel time?


